\section{Dataset Schema and Views}
\label{sec:schema-views}

Each released row corresponds to exactly one candidate victim at exactly
one full-cache-miss decision: which trace and split it came from, the
cache configuration (capacity, horizon), which decision it belongs to,
the candidate's identity, its label (\texttt{y\_loss}, the finite-horizon
counterfactual miss cost defined in Section~\ref{sec:label-generation},
and its sign-flipped alias \texttt{y\_value}), and a set of request/
candidate/cache features. Table~\ref{tab:schema-summary} lists the full
column groups; this section describes only the two derived views built on
top of these rows.

Grouping candidate rows by decision identifier gives the \emph{decision
view}: candidate count, minimum loss, the set of optimal candidates, tie
count, and regret summary statistics for that decision (this is exactly
the aggregation illustrated by hand in Section~\ref{sec:background}).
Comparing candidate pairs within a decision and ordering them by
\texttt{y\_loss} gives the \emph{pairwise view}, with ties when losses are
equal; the evaluated open release ships a capped 1,000,000-row pairwise
sample rather than the full pairwise view. The three views are layered
projections of the same underlying decisions -- candidate rows are the
atomic supervised objects, the decision view groups them, and the pairwise
sample keeps selected within-decision comparisons -- not three separate
sources of ground truth.

\begin{table}[t]
  \caption{Compact summary of candidate-row schema groups and derived-view fields.}
  \label{tab:schema-summary}
  \centering
  \footnotesize
  \begin{tabularx}{\columnwidth}{@{}P{0.31\columnwidth}L@{}}
    \toprule
    Column group & Notes \\
    \midrule
    Decision metadata & family, capacity, horizon, split, IDs \\
    Candidate identity & candidate page ID \\
    Labels & \texttt{y\_loss}, \texttt{y\_value} \\
    Feature groups & request, candidate, cache, disagreement, rate \\
    Derived decision fields & tie and regret summaries \\
    Derived pairwise fields & A/B labels and regret differences \\
    \bottomrule
  \end{tabularx}
\end{table}
